\section{Fundamentação teórica}
\label{sec:fundamentacao}

\subsection{Erro humano e barreiras de prevenção}

\cite{reason1990} diferencia formas de falha associadas à execução e ao
planejamento e mostra que o desempenho humano é influenciado pelo sistema em
que a tarefa ocorre. Essa perspectiva evita reduzir o problema a uma culpa
individual e direciona a análise para sinais, sequência da tarefa, interface,
equipamento e condições de trabalho. No abastecimento, \cite{adamsDavid2007}
aplicaram análise hierárquica da tarefa e encontraram fatores físicos e
psicológicos que, isolados ou combinados, abrem caminhos para a seleção errada.

O intertravamento proposto é uma barreira de prevenção: a ação física depende de
uma autorização calculada antes da retirada do bico de seu suporte. A barreira
atua sobre o acesso ao bico, não sobre o bombeamento ou a vazão. Isso não elimina
treinamento, conferência e procedimentos. Também pode criar novos modos de
falha, como indisponibilidade do serviço, código veicular não reconhecido,
comando atrasado ou uso abusivo de uma exceção. Por isso, o projeto adota os
princípios de autorização explícita, menor privilégio e padrão seguro diante de
falhas, coerentes com princípios clássicos de proteção de sistemas
\cite{saltzerSchroeder1975}.

\subsection{Consulta veicular e política de compatibilidade}

O serviço Consulta Online SENATRAN expõe operações de consulta a bases mantidas pela
Secretaria Nacional de Trânsito. No inventário público, os campos ``Código
Combustível'' e ``Combustível Descrição'' aparecem na consulta detalhada de
veículo \cite{serproDadosVeiculo2026}. O retorno do provedor, contudo, não deve
acionar diretamente um dispositivo. Entre a consulta e o atuador haverá uma
camada de política que:

\begin{enumerate}
  \item valida formato, atualidade e completude da resposta;
  \item converte códigos externos para categorias internas versionadas;
  \item representa veículos com um ou mais combustíveis admitidos;
  \item compara o produto solicitado à lista de compatibilidade; e
  \item produz uma decisão auditável: autorizado, negado ou revisão necessária.
\end{enumerate}

Essa separação permite substituir o provedor real por dados sintéticos durante
o desenvolvimento e testar códigos novos sem expor dados pessoais. Também evita
acoplar a segurança física a detalhes instáveis da API.

\subsection{Sistemas ciberfísicos, mensageria e estado seguro}

O protótipo combina software de negócio, comunicação e atuadores e foi concebido
para que uma decisão digital possa comandar um efeito físico observável. A
validação do movimento mecânico, porém, permanece pendente. A comunicação entre o backend
e o controlador será representada por mensagens MQTT, protocolo padronizado pela
OASIS com níveis de qualidade de serviço e sessões próprias
\cite{oasisMqtt2019}. Esses recursos não garantem, sozinhos, que um comando seja
executado uma única vez. O projeto acrescenta identificador idempotente,
validade temporal, confirmação e bloqueio local por expiração.

O estado seguro comandado do protótipo é manter desenergizadas as saídas de
abertura de todas as travas. Inicialização, perda de conexão, resposta ambígua,
comando vencido e término da sessão conduzem a esse estado. Se um bico já
estiver fora do suporte, desenergizar a saída não o recoloca fisicamente; sua
devolução continua sendo uma ação humana. Uma liberação excepcional não altera
essa premissa: ela exige supervisor autenticado, justificativa e trilha de
auditoria.

\subsection{Proteção de dados e acesso regulado}

A placa pode viabilizar associação a uma pessoa natural no contexto do
tratamento e, por prudência, será tratada como dado pessoal. A Lei Geral de
Proteção de Dados Pessoais estabelece, entre outros, os princípios de finalidade,
adequação, necessidade, segurança, prevenção e responsabilização
\cite{lgpd2018}. A Portaria SENATRAN n\textsuperscript{o} 139/2025 disciplina o acesso aos sistemas
e subsistemas, incluindo avaliação do caso de uso, controles de segurança e
atuação do Serpro e de uma GCC \cite{senatran1392025}.

Consequentemente, a versão de bancada utiliza placas e respostas sintéticas.
Na integração autorizada, a placa será usada apenas durante a decisão e não
constará do log padrão. Também não serão persistidos resposta bruta, nome de
proprietário, Renavam ou chassi. A base legal, os avisos ao titular, a retenção e
as responsabilidades deverão refletir o caso de uso efetivamente aprovado, não
uma suposição do projeto.

\subsection{Segurança da instalação}

A NR-20 define requisitos de segurança e saúde no trabalho com inflamáveis e
combustíveis e contempla o conceito de área classificada \cite{nr20}. A Portaria
Inmetro n\textsuperscript{o} 115/2022 trata da avaliação da conformidade de equipamentos elétricos
para atmosferas explosivas \cite{inmetro1152022}. Assim, ESP32, servo, solenoide
e tablet comuns não serão instalados na zona de abastecimento. O ensaio acadêmico
ocorrerá sem líquido inflamável, e uma eventual evolução de campo dependerá de
projeto específico do intertravamento independente, profissionais habilitados e
equipamentos certificados para a classificação do local. Essa evolução não
pressupõe alteração da bomba ou integração com a automação da empresa.
